\chapter[Lezione 3 --- 1 ottobre 2026]{Lezione 3: coefficienti binomiali e numeri reali}
\label{chap:lezione-03}

\section{Coefficienti binomiali}
\label{sec:lez03-coefficienti-binomiali}

La definizione del fattoriale introdotta nella \cref{def:fattoriale}
permette di definire i coefficienti binomiali. Essi saranno considerati
per ogni coppia di numeri naturali.

\begin{definition}[Coefficiente binomiale]
\label{def:coefficiente-binomiale}
Siano \(n,k\in\N\). Si definisce
\[
  \binom{n}{k}
  :=
  \begin{cases}
    \displaystyle\frac{n!}{k!(n-k)!}, & \text{se } k\le n,\\[8pt]
    0, & \text{se } k>n.
  \end{cases}
\]
\end{definition}

Dalla definizione si ricavano immediatamente alcuni casi particolari. Per
ogni \(n\in\N\),
\[
  \binom{n}{0}
  =\frac{n!}{0!\,n!}
  =1
\]
e
\[
  \binom{n}{n}
  =\frac{n!}{n!\,0!}
  =1.
\]
Inoltre, per ogni \(n\in\N\) con \(n\ge1\),
\begin{align*}
  \binom{n}{n-1}
    &=\frac{n!}{(n-1)!\,(n-(n-1))!}\\
    &=\frac{n!}{(n-1)!\,1!}\\
    &=n.
\end{align*}

Una prima proprietà fondamentale dei coefficienti binomiali è la loro
simmetria rispetto agli indici \(k\) e \(n-k\).

\begin{proposition}[Simmetria dei coefficienti binomiali]
\label{prop:simmetria-coefficienti-binomiali}
Per ogni \(n,k\in\N\) con \(k\le n\) vale
\[
  \binom{n}{k}=\binom{n}{n-k}.
\]
\end{proposition}

\begin{proof}
Per la \cref{def:coefficiente-binomiale},
\[
  \binom{n}{k}
  =\frac{n!}{k!(n-k)!}.
\]
D'altra parte,
\begin{align*}
  \binom{n}{n-k}
    &=\frac{n!}{(n-k)!\,(n-(n-k))!}\\
    &=\frac{n!}{(n-k)!\,k!}.
\end{align*}
Poiché il prodotto è commutativo, i due denominatori coincidono e quindi
\[
  \binom{n}{k}=\binom{n}{n-k}.
\]
\end{proof}

Un'altra relazione fondamentale è l'identità di Pascal, che esprime un
coefficiente binomiale come somma di due coefficienti della riga
precedente.

\begin{proposition}[Identità di Pascal]
\label{prop:identita-pascal}
Per ogni \(n,k\in\N\) con \(n\ge1\) e \(1\le k\le n\) vale
\[
  \binom{n}{k}
  =\binom{n-1}{k-1}+\binom{n-1}{k}.
\]
\end{proposition}

\begin{proof}
Si considera anzitutto il caso \(1\le k\le n-1\). Applicando la
\cref{def:coefficiente-binomiale}, si ottiene
\begin{align*}
  \binom{n-1}{k-1}+\binom{n-1}{k}
    &=\frac{(n-1)!}{(k-1)!(n-k)!}
      +\frac{(n-1)!}{k!(n-k-1)!}.
\end{align*}

Poiché
\[
  (n-k)!=(n-k)(n-k-1)!
  \qquad\text{e}\qquad
  k!=k(k-1)!,
\]
si può raccogliere il fattore comune
\[
  \frac{(n-1)!}{(k-1)!(n-k-1)!}
\]
e ottenere
\begin{align*}
  \binom{n-1}{k-1}+\binom{n-1}{k}
    &=\frac{(n-1)!}{(k-1)!(n-k-1)!}
      \left(\frac{1}{n-k}+\frac{1}{k}\right)\\
    &=\frac{(n-1)!}{(k-1)!(n-k-1)!}
      \frac{n}{k(n-k)}\\
    &=\frac{n!}{k!(n-k)!}\\
    &=\binom{n}{k}.
\end{align*}

Resta il caso \(k=n\). Per la definizione estesa dei coefficienti
binomiali,
\[
  \binom{n-1}{n}=0,
\]
mentre
\[
  \binom{n-1}{n-1}=1
  \qquad\text{e}\qquad
  \binom{n}{n}=1.
\]
Pertanto anche in questo caso
\[
  \binom{n}{n}
  =\binom{n-1}{n-1}+\binom{n-1}{n}.
\]
\end{proof}

Benché la formula della \cref{def:coefficiente-binomiale} sia espressa
come rapporto di fattoriali, il suo valore è sempre un numero naturale:
\[
  \binom{n}{k}\in\N
  \qquad
  \forall n,k\in\N.
\]
La dimostrazione di questa proprietà non viene sviluppata in questa sede.

\clearpage

\section{Numeri reali}
\label{sec:lez03-numeri-reali}

Gli insiemi numerici introdotti finora sono legati dalla catena di
inclusioni
\[
  \N\subseteq\Z\subseteq\Q.
\]
Come mostrato nel \cref{thm:irrazionalita-radice-due}, l'insieme dei
razionali non è tuttavia sufficiente a contenere tutte le soluzioni di
equazioni algebriche elementari. Si introduce pertanto un insieme numerico
più ampio, l'insieme dei numeri reali \(\R\), ottenendo
\[
  \N\subseteq\Z\subseteq\Q\subseteq\R.
\]

Nel seguito si utilizza una proprietà di completezza di \(\R\), formulata
mediante il cosiddetto assioma di separazione.

Siano \(A,B\subseteq\R\) due insiemi non vuoti. Si dice che \(A\) e \(B\)
sono \emph{separati}, con \(A\) a sinistra di \(B\), se
\[
  a\le b
  \qquad
  \forall a\in A,\ \forall b\in B.
\]

\phantomsection
\label{ass:assioma-separazione}
\textbf{Assioma di separazione.}
Se \(A,B\subseteq\R\) sono non vuoti e separati, allora esiste almeno un
numero reale \(x\) tale che
\[
  a\le x\le b
  \qquad
  \forall a\in A,\ \forall b\in B.
\]
Un tale numero \(x\) è detto \emph{elemento separatore} degli insiemi
\(A\) e \(B\).

L'assioma di separazione esprime una proprietà caratteristica dei numeri
reali. Esistono altre formulazioni equivalenti della completezza di
\(\R\), ma non vengono sviluppate in questa sede.

\section{Proprietà archimedea}
\label{sec:lez03-proprieta-archimedea}

L'assioma di separazione permette di dimostrare che i numeri naturali non
sono superiormente limitati nell'insieme dei reali.

\begin{theorem}[Proprietà archimedea]
\label{thm:proprieta-archimedea}
Per ogni \(x\in\R\) esiste \(n\in\N\) tale che
\[
  n>x.
\]
Equivalentemente,
\[
  \forall x\in\R\;\exists n\in\N:\ n>x.
\]
\end{theorem}

\begin{proof}
Sia \(x\in\R\).

Se \(x\le0\), è sufficiente scegliere \(n=1\), poiché
\[
  1>x.
\]

Si considera quindi il caso \(x>0\) e si pone
\[
  A=\{n\in\N:n\le x\}.
\]
Poiché \(0\le x\), si ha \(0\in A\), dunque \(A\neq\varnothing\).
L'obiettivo è mostrare che \(A\neq\N\). Infatti, da questa disuguaglianza
seguirebbe l'esistenza di un naturale \(n\notin A\), cioè di un naturale
per cui \(n>x\).

Si procede per assurdo supponendo
\[
  A=\N.
\]
Si definisce allora
\[
  B=\{y\in\R:y\ge n\ \text{per ogni }n\in\N\}.
\]
Dall'ipotesi \(A=\N\) segue che
\[
  n\le x
  \qquad
  \forall n\in\N,
\]
perciò \(x\in B\) e quindi \(B\neq\varnothing\).

Inoltre, per ogni \(n\in A=\N\) e per ogni \(y\in B\), dalla definizione
di \(B\) si ha
\[
  n\le y.
\]
Gli insiemi \(A\) e \(B\) sono dunque separati. Per l'\hyperref[ass:assioma-separazione]{assioma di separazione} esiste
\(\lambda\in\R\) tale che
\[
  n\le\lambda\le y
  \qquad
  \forall n\in A,\ \forall y\in B.
\]

Poiché \(A=\N\), per ogni \(n\in\N\) anche \(n+1\in A\). Dalla proprietà
di separazione segue quindi
\[
  n+1\le\lambda,
\]
da cui
\[
  n\le\lambda-1
  \qquad
  \forall n\in\N.
\]
Per definizione di \(B\), ciò implica
\[
  \lambda-1\in B.
\]

D'altra parte, essendo \(\lambda\) un elemento separatore, vale
\[
  \lambda\le y
  \qquad
  \forall y\in B.
\]
Scegliendo in particolare \(y=\lambda-1\), si ottiene
\[
  \lambda\le\lambda-1,
\]
ovvero
\[
  1\le0,
\]
che è assurdo.

L'ipotesi \(A=\N\) è pertanto falsa. Esiste quindi almeno un
\(n\in\N\setminus A\), e dalla definizione di \(A\) segue
\[
  n>x.
\]
\end{proof}

\clearpage

\section{Parte intera, floor e ceiling}
\label{sec:lez03-parte-intera}

La proprietà archimedea conduce a un'importante caratterizzazione della
posizione di un numero reale rispetto agli interi consecutivi.

\begin{proposition}[Proprietà della parte intera]
\label{prop:parte-intera}
Per ogni \(x\in\R\) esiste un unico \(n\in\Z\) tale che
\[
  n\le x<n+1.
\]
In simboli,
\[
  \forall x\in\R\;\exists!\,n\in\Z:\ n\le x<n+1.
\]
\end{proposition}

La \cref{prop:parte-intera} garantisce quindi che ogni numero reale
\(x\) individua un unico intero \(n\) che non supera \(x\), mentre
l'intero successivo \(n+1\) lo supera strettamente. Tale intero coincide
con il più grande intero minore o uguale a \(x\) e viene chiamato
\emph{parte intera} di \(x\).

Questa proprietà permette di introdurre una funzione che associa a ogni
numero reale la sua parte intera.

\begin{definition}[Funzione floor]
\label{def:floor}
Per ogni \(x\in\R\), si definisce
\[
  \lfloor x\rfloor=n,
\]
dove \(n\in\Z\) è l'unico intero individuato dalla
\cref{prop:parte-intera}.

Equivalentemente,
\[
  \lfloor x\rfloor\le x<\lfloor x\rfloor+1.
\]
\end{definition}

La funzione floor associa quindi a ogni numero reale il più grande intero
che non lo supera. Per esempio,
\[
  \lfloor x\rfloor=0
  \qquad
  \text{per ogni }x\in[0,1).
\]
Più in generale, per ogni \(n\in\Z\),
\[
  \lfloor x\rfloor=n
  \quad\Longleftrightarrow\quad
  n\le x<n+1
  \quad\Longleftrightarrow\quad
  x\in[n,n+1).
\]

\begin{figure}[htbp]
  \centering
  \begin{minipage}[c]{0.56\textwidth}
    \centering
    \begin{tikzpicture}[unipd/diagram,scale=0.9]
  \fill[green!12] (0,-0.3) rectangle (1,0.3);
  \draw[unipd-rule,very thin] (-2,-2) grid (2,2);
  \draw[->,unipd-muted] (-2.4,0) -- (2.6,0) node[right] {\(x\)};
  \draw[->,unipd-muted] (0,-2.4) -- (0,2.6) node[above] {\(y\)};

  \foreach \k in {-2,-1,0,1,2} {
    \draw[unipd-muted] (\k,-2.45) -- (\k,-2.55)
      node[below] {\(\k\)};
    \draw[unipd-muted] (-2.45,\k) -- (-2.55,\k)
      node[left] {\(\k\)};
  }

  \foreach \n in {-2,-1,0,1} {
    \draw[very thick,unipd-primary] (\n,\n) -- ({\n+1},\n);
    \fill[unipd-primary] (\n,\n) circle (2pt);
    \draw[unipd-primary,fill=white,thick] ({\n+1},\n) circle (2pt);
  }
  \draw[very thick,unipd-primary] (2,2) -- (2.4,2);
  \fill[unipd-primary] (2,2) circle (2pt);

  \draw[very thick,green!40!black] (0,0) -- (1,0);
  \fill[green!40!black] (0,0) circle (2pt);
  \draw[green!40!black,fill=white,thick] (1,0) circle (2pt);
  \node[fill=green!12,draw=green!40!black,rounded corners,inner sep=3pt]
    at (0.5,-0.55) {\(\lfloor x\rfloor=0\), \(x\in[0,1)\)};
\end{tikzpicture}

  \end{minipage}
  \hfill
  \begin{minipage}[c]{0.38\textwidth}
    \small
    \textbf{Esempi in Python}
    \begin{unipdcode}[language=Python,numbers=none,xleftmargin=0pt,framexleftmargin=0pt]
from math import floor

print(floor(0.5))   # Output: 0
print(floor(-0.5))  # Output: -1
print(floor(1.0))   # Output: 1
\end{unipdcode}

    \vspace{0.5em}
  \end{minipage}
  \caption{Grafico della funzione floor ed esempi d'uso in Python.}
  \label{fig:floor}
\end{figure}

In modo analogo si introduce la funzione ceiling, che associa a ogni
numero reale il più piccolo intero maggiore o uguale al numero stesso.

\begin{definition}[Funzione ceiling]
\label{def:ceiling}
Per ogni \(x\in\R\), si definisce
\[
  \lceil x\rceil=m,
\]
dove \(m\in\Z\) è l'unico intero tale che
\[
  m-1<x\le m.
\]
Equivalentemente,
\[
  \lceil x\rceil-1<x\le\lceil x\rceil.
\]
\end{definition}

La funzione ceiling associa quindi a ogni numero reale il più piccolo intero
maggiore o uguale al numero stesso. Per esempio,
\[
  \lceil x\rceil=1
  \qquad
  \text{per ogni }x\in(0,1].
\]
Più in generale, per ogni \(m\in\Z\),
\[
  \lceil x\rceil=m
  \quad\Longleftrightarrow\quad
  m-1<x\le m
  \quad\Longleftrightarrow\quad
  x\in(m-1,m].
\]

\begin{figure}[htbp]
  \centering
  \begin{minipage}[c]{0.56\textwidth}
    \centering
    \begin{tikzpicture}[unipd/diagram,scale=0.9]
  \fill[green!12] (0,0.7) rectangle (1,1.3);
  \draw[unipd-rule,very thin] (-2,-2) grid (2,2);
  \draw[->,unipd-muted] (-2.4,0) -- (2.6,0) node[right] {\(x\)};
  \draw[->,unipd-muted] (0,-2.4) -- (0,2.6) node[above] {\(y\)};

  \foreach \k in {-2,-1,0,1,2} {
    \draw[unipd-muted] (\k,-2.45) -- (\k,-2.55)
      node[below] {\(\k\)};
    \draw[unipd-muted] (-2.45,\k) -- (-2.55,\k)
      node[left] {\(\k\)};
  }

  \draw[very thick,unipd-primary] (-2.4,-2) -- (-2,-2);
  \fill[unipd-primary] (-2,-2) circle (2pt);
  \foreach \m in {-1,0,1,2} {
    \draw[very thick,unipd-primary] ({\m-1},\m) -- (\m,\m);
    \draw[unipd-primary,fill=white,thick] ({\m-1},\m) circle (2pt);
    \fill[unipd-primary] (\m,\m) circle (2pt);
  }

  \draw[very thick,green!40!black] (0,1) -- (1,1);
  \draw[green!40!black,fill=white,thick] (0,1) circle (2pt);
  \fill[green!40!black] (1,1) circle (2pt);
  \node[fill=green!12,draw=green!40!black,rounded corners,inner sep=3pt]
    at (0.5,-0.55) {\(\lceil x\rceil=1\), \(x\in(0,1]\)};
\end{tikzpicture}

  \end{minipage}
  \hfill
  \begin{minipage}[c]{0.38\textwidth}
    \small
    \textbf{Esempi in Python}
    \begin{unipdcode}[language=Python,numbers=none,xleftmargin=0pt,framexleftmargin=0pt]
from math import ceil

print(ceil(0.5))   # Output: 1
print(ceil(-0.5))  # Output: 0
print(ceil(1.0))   # Output: 1
\end{unipdcode}

    \vspace{0.5em}
  \end{minipage}
  \caption{Grafico della funzione ceiling ed esempi d'uso in Python.}
  \label{fig:ceiling}
\end{figure}

\section[Cenno sulla radice quadrata di 2 nei reali]
{Cenno su \(\sqrt{2}\in\R\setminus\Q\)}
\label{sec:lez03-completezza-radice-due}

Nella \cref{sec:lez02-irrazionalita-radice-due} è stato dimostrato che
\[
  \sqrt{2}\notin\Q.
\]

Per motivare l'esistenza di \(\sqrt{2}\) nell'insieme dei numeri reali,
si considerano gli insiemi
\[
  A=\{x\in\R:x\ge0,\ x^2<2\},
  \qquad
  B=\{x\in\R:x\ge0,\ x^2>2\}.
\]

Resta da verificare che \(A\) e \(B\) siano separati. L'idea è quindi
applicare l'assioma di separazione per individuare un elemento
separatore \(\lambda\in\R\). Geometricamente, tale elemento \(\lambda\) si colloca tra i due
insiemi
\[
  a\le\lambda\le b
  \qquad
  \forall a\in A,\ \forall b\in B.
\]

Il passaggio conclusivo consiste nel mostrare che
\[
  \lambda^2=2.
\]

Il cenno conduce quindi alla conclusione \(\sqrt{2}\in\R\setminus\Q\).